% GENERATED FILE. Edit canonical lessons, then run npm run generate:books.
% canonical-source-hash: fnv1a64:e8ed1845f93f31c6
% canonical-lessons: GU-C381-pastavu, GU-C381-aphsos, GU-C381-kas, GU-C381-jhurvu, GU-C381-moko

\chapter{If Only}
\label{ch:gu-if-only}
\hlchaptermodality{381}

\begin{chapteropening}
\textbf{By the end of this chapter:} \emph{I can say what I wish had happened.}

Complete the last lesson of chapter 381: I can say what I wish had happened.
\end{chapteropening}

\section[\texorpdfstring{\gu{પસ્તાવું} (pastāvũ)}{પસ્તાવું (pastāvũ)}]{\gu{પસ્તાવું} (pastāvũ) — to regret}

\label{lesson:GU-C381-pastavu}



\begin{quote}
Before the new one: say the Gujarati for slightly, then the Gujarati for to cope.
\end{quote}

\subsection*{You'll want to know: \gu{પસ્તાવું}}
\textbf{\gu{પસ્તાવું}} — \emph{pastāvũ} — \textquotedblleft{}to regret\textquotedblright{}.

\begin{quote}
\gu{જો} \gu{તમે} \gu{આ} \gu{તક} \gu{ખોશો} \gu{તો} \gu{પછી} \gu{પસ્તાશો}.
\end{quote}

\begin{quote}
\emph{jo tame ā tak khośo to pachī pastāśo.}
\end{quote}

\begin{quote}
If you lose this chance, you will regret it later.
\end{quote}

Past (he …): \gu{તે} \gu{પોતાની} \gu{ભૂલ} \gu{પર} \gu{પસ્તાયો}, \emph{te potānī bhūl par pastāyo}.

\begin{cousinweb}
\textbf{\gu{પસ્તાવું}} (\emph{pastāvũ}) CDIAL lists the Gujarati word under a reconstructed \emph{paścōttāpa}, \textquotedblleft{}repentance\textquotedblright{} (CDIAL 8010), which is built from \emph{paścāt}, \textquotedblleft{}after\textquotedblright{}, and \emph{tāpa}, \textquotedblleft{}burning\textquotedblright{}: the heat that comes afterwards.
\end{cousinweb}

\begin{grammarlens}[title={What you've built}]
One more word for saying what happens if.
\end{grammarlens}

\subsection*{Your turn}
\begin{itemize}
\raggedright
  \item \emph{Say it:} \emph{pastāvũ}
  \item \emph{Say it:} \emph{pastāvũ}, once more
  \item \emph{Say it:} \emph{jo tame ā tak khośo to pachī pastāśo.} — If you lose this chance, you will regret it later.
  \item \emph{Recall:} say the Gujarati for slightly, then the Gujarati for to cope, then say \emph{pastāvũ} again
\end{itemize}

\begin{tcolorbox}[breakable,colback=teal!4,colframe=teal!35!black,title={Before you move on}]
What does \gu{પસ્તાવું} mean? (\textquotedblleft{}To regret\textquotedblright{}.) Say it once more.
\end{tcolorbox}
\section[\texorpdfstring{\gu{અફસોસ} (aphsos)}{અફસોસ (aphsos)}]{\gu{અફસોસ} (aphsos) — a pity; regret}

\label{lesson:GU-C381-aphsos}



\begin{quote}
Before the new one: say the Gujarati for to cope, then the Gujarati for to regret.
\end{quote}

\subsection*{You'll want to know: \gu{અફસોસ}}
\textbf{\gu{અફસોસ}} — \emph{aphsos} — \textquotedblleft{}a pity; regret\textquotedblright{}.

\begin{quote}
\gu{જો} \gu{હું} \gu{ત્યાં} \gu{હોત} \gu{તો} \gu{તમને} \gu{મદદ} \gu{કરત}; \gu{મને} \gu{તેનો} \gu{અફસોસ} \gu{છે}.
\end{quote}

\begin{quote}
\emph{jo hũ tyā̃ hot to tamne madad karat; mane teno aphsos che.}
\end{quote}

\begin{quote}
If I had been there I would have helped you; I am sorry for that.
\end{quote}

A masculine noun.

\begin{cousinweb}
\textbf{\gu{અફસોસ}} (\emph{aphsos}) is usually traced to Persian \emph{afsōs}, \textquotedblleft{}regret, alas\textquotedblright{}.
\end{cousinweb}

\begin{grammarlens}[title={What you've built}]
One more word for saying what happens if.
\end{grammarlens}

\subsection*{Your turn}
\begin{itemize}
\raggedright
  \item \emph{Say it:} \emph{aphsos}
  \item \emph{Say it:} \emph{aphsos}, once more
  \item \emph{Say it:} \emph{jo hũ tyā̃ hot to tamne madad karat; mane teno aphsos che.} — If I had been there I would have helped you; I am sorry for that.
  \item \emph{Recall:} say the Gujarati for to cope, then the Gujarati for to regret, then say \emph{aphsos} again
\end{itemize}

\begin{tcolorbox}[breakable,colback=teal!4,colframe=teal!35!black,title={Before you move on}]
What does \gu{અફસોસ} mean? (\textquotedblleft{}A pity; regret\textquotedblright{}.) Say it once more.
\end{tcolorbox}
\section[\texorpdfstring{\gu{કાશ} (kāś)}{કાશ (kāś)}]{\gu{કાશ} (kāś) — if only}

\label{lesson:GU-C381-kas}



\begin{quote}
Before the new one: say the Gujarati for to regret, then the Gujarati for a pity.
\end{quote}

\subsection*{You'll want to know: \gu{કાશ}}
\textbf{\gu{કાશ}} — \emph{kāś} — \textquotedblleft{}if only\textquotedblright{}.

\begin{quote}
\gu{કાશ} \gu{મારી} \gu{પાસે} \gu{વધારે} \gu{સમય} \gu{હોત}!
\end{quote}

\begin{quote}
\emph{kāś mārī pāse vadhāre samay hot!}
\end{quote}

\begin{quote}
If only I had more time!
\end{quote}

\begin{cousinweb}
\textbf{\gu{કાશ}} (\emph{kāś}) is usually traced to Persian \emph{kāsh}, \textquotedblleft{}would that\textquotedblright{}; in Gujarati it opens a wish that did not come true.
\end{cousinweb}

\begin{grammarlens}[title={What you've built}]
One more word for saying what happens if.
\end{grammarlens}

\subsection*{Your turn}
\begin{itemize}
\raggedright
  \item \emph{Say it:} \emph{kāś}
  \item \emph{Say it:} \emph{kāś}, once more
  \item \emph{Say it:} \emph{kāś mārī pāse vadhāre samay hot!} — If only I had more time!
  \item \emph{Recall:} say the Gujarati for to regret, then the Gujarati for a pity, then say \emph{kāś} again
\end{itemize}

\begin{tcolorbox}[breakable,colback=teal!4,colframe=teal!35!black,title={Before you move on}]
What does \gu{કાશ} mean? (\textquotedblleft{}If only\textquotedblright{}.) Say it once more.
\end{tcolorbox}
\section[\texorpdfstring{\gu{ઝૂરવું} (jhūrvũ)}{ઝૂરવું (jhūrvũ)}]{\gu{ઝૂરવું} (jhūrvũ) — to long for someone; to pine}

\label{lesson:GU-C381-jhurvu}



\begin{quote}
Before the new one: say the Gujarati for a pity, then the Gujarati for if only.
\end{quote}

\subsection*{You'll want to know: \gu{ઝૂરવું}}
\textbf{\gu{ઝૂરવું}} — \emph{jhūrvũ} — \textquotedblleft{}to long for someone; to pine\textquotedblright{}.

\begin{quote}
\gu{જો} \gu{દીકરો} \gu{શહેરમાં} \gu{રહેવા} \gu{જશે} \gu{તો} \gu{માતા} \gu{તેના} \gu{માટે} \gu{ઝૂરશે}.
\end{quote}

\begin{quote}
\emph{jo dīkro śahermā̃ rahevā jaśe to mātā tenā māṭe jhūrśe.}
\end{quote}

\begin{quote}
If her son goes to live in the city, the mother will pine for him.
\end{quote}

Past (he …): \gu{તે} \gu{ઘર} \gu{માટે} \gu{ઘણું} \gu{ઝૂર્યો}, \emph{te ghar māṭe ghaṇũ jhūryo}.

\begin{cousinweb}
\textbf{\gu{ઝૂરવું}} (\emph{jhūrvũ}) CDIAL lists the Gujarati verb, \textquotedblleft{}to languish\textquotedblright{}, under a reconstructed \emph{jhūrati}, \textquotedblleft{}wastes away\textquotedblright{} (CDIAL 5409).
\end{cousinweb}

\begin{grammarlens}[title={What you've built}]
One more word for saying what happens if.
\end{grammarlens}

\subsection*{Your turn}
\begin{itemize}
\raggedright
  \item \emph{Say it:} \emph{jhūrvũ}
  \item \emph{Say it:} \emph{jhūrvũ}, once more
  \item \emph{Say it:} \emph{jo dīkro śahermā̃ rahevā jaśe to mātā tenā māṭe jhūrśe.} — If her son goes to live in the city, the mother will pine for him.
  \item \emph{Recall:} say the Gujarati for a pity, then the Gujarati for if only, then say \emph{jhūrvũ} again
\end{itemize}

\begin{tcolorbox}[breakable,colback=teal!4,colframe=teal!35!black,title={Before you move on}]
What does \gu{ઝૂરવું} mean? (\textquotedblleft{}To long for someone; to pine\textquotedblright{}.) Say it once more.
\end{tcolorbox}
\section[\texorpdfstring{\gu{મોકો} (moko)}{મોકો (moko)}]{\gu{મોકો} (moko) — an opening; a chance}

\label{lesson:GU-C381-moko}



\begin{quote}
Before the new one: say the Gujarati for if only, then the Gujarati for to long for someone.
\end{quote}

\subsection*{You'll want to know: \gu{મોકો}}
\textbf{\gu{મોકો}} — \emph{moko} — \textquotedblleft{}an opening; a chance\textquotedblright{}.

\begin{quote}
\gu{જો} \gu{મારી} \gu{પાસે} \gu{આ} \gu{મોકો} \gu{હોત} \gu{તો} \gu{હું} \gu{તરત} \gu{તે} \gu{લેત}.
\end{quote}

\begin{quote}
\emph{jo mārī pāse ā moko hot to hũ tarat te let.}
\end{quote}

\begin{quote}
If I had this opportunity, I would take it at once.
\end{quote}

A masculine noun.

\begin{cousinweb}
\textbf{\gu{મોકો}} (\emph{moko}) is usually traced to Arabic \emph{mauqaʿ}, \textquotedblleft{}a place, an occasion\textquotedblright{}, through Persian.
\end{cousinweb}

\begin{grammarlens}[title={What you've built}]
One more word for saying what happens if.
\end{grammarlens}

\subsection*{Your turn}
\begin{itemize}
\raggedright
  \item \emph{Say it:} \emph{moko}
  \item \emph{Say it:} \emph{moko}, once more
  \item \emph{Say it:} \emph{jo mārī pāse ā moko hot to hũ tarat te let.} — If I had this opportunity, I would take it at once.
  \item \emph{Recall:} say the Gujarati for if only, then the Gujarati for to long for someone, then say \emph{moko} again
\end{itemize}

\begin{tcolorbox}[breakable,colback=teal!4,colframe=teal!35!black,title={Before you move on}]
What does \gu{મોકો} mean? (\textquotedblleft{}An opening; a chance\textquotedblright{}.) Say it once more.
\end{tcolorbox}
